\section{Il lavoro svolto}
\label{sec:lavoro}

\subsection{Approccio generale}
\label{sec:approccio}

Il lavoro mette insieme due programmi che fanno cose diverse.
OpenFOAM~\cite{openfoam13} è la piattaforma CFD: gestisce la mesh, risolve le
equazioni di conservazione cella per cella e fa avanzare la soluzione nel
tempo. \mpp~\cite{scoggins2020} è una libreria di termochimica: dato lo stato
di un gas, calcola quanta energia contiene ciascun modo, quanto dura il
rilassamento vibrazionale e quanto sono veloci le reazioni, e sa fare il
percorso inverso, dalle energie alle temperature. L'idea è lasciare a ciascuno
il proprio mestiere: OpenFOAM risolve le equazioni del modello a due
temperature, \mpp{} fornisce tutte le grandezze termochimiche che vi
compaiono. Il solver, in particolare, non segue direttamente le due
temperature ma le energie corrispondenti, che si conservano correttamente anche
attraverso un urto, e a ogni passo le consegna a \mpp, che ne ricava le
temperature.

Il collegamento con OpenFOAM è realizzato senza modificarne il codice,
attraverso \emph{librerie dinamiche}. Una libreria dinamica è un pezzo di codice
compilato che non fa parte di un programma, ma viene caricato da questo mentre
è in esecuzione, come un'estensione aggiunta a un'applicazione già installata.
OpenFOAM è costruito attorno a questo meccanismo: il solver, il modello
termodinamico, le condizioni al contorno e molti altri componenti non sono
fissati nel programma, ma scelti per nome nei file di testo del caso. Ogni
libreria, quando viene caricata, iscrive i componenti che contiene in un
elenco interno di OpenFOAM; OpenFOAM legge nel caso il nome richiesto, lo cerca
nell'elenco e costruisce il componente corrispondente. Aggiungere un modello
nuovo significa quindi compilare una libreria che lo iscriva nell'elenco
giusto, e OpenFOAM lo userà come uno dei propri.

\subsection{OpenFOAM}
\label{sec:openfoam}

\subsubsection{Com'è fatto e il suo modello a una temperatura}
\label{sec:openfoam-moduli}

OpenFOAM è una piattaforma CFD a codice aperto: un insieme di librerie e di
solver già pronti che discretizzano le equazioni con il metodo dei volumi
finiti, dividendo il dominio in celle e tenendo un valore di ogni grandezza per
cella. È scritta in modo che si possa aggiungere fisica nuova senza toccarne il
nucleo, ed è proprio quello che fa il progetto.

Per i flussi che interessano qui fa già molto. Uno dei suoi solver è fatto
apposta per i flussi compressibili veloci: gestisce le onde d'urto, trasporta
le specie chimiche insieme al flusso e risolve un'equazione per l'energia. A chiudere il
sistema c'è un modello termodinamico: dall'energia di ogni cella ricava una
temperatura, e da quella la pressione, la viscosità, la conducibilità e le
velocità di reazione. Tutta la catena, però, poggia su una sola temperatura: è
il modello a una temperatura della sezione~\ref{sec:una-temperatura}.

\subsubsection{Dove si innesta il progetto}
\label{sec:openfoam-innesto}

Il progetto si innesta in due punti di questa catena. Il suo solver deriva da quel solver per flussi
compressibili veloci, lo stesso da cui è nato \hyfoam, e ne riscrive solo la fase dell'energia: lì aggiunge
l'equazione dell'energia vibro-elettronica e i termini sorgente della chimica,
mentre flussi, densità, quantità di moto e pressione restano quelli ereditati.
Il suo modello termodinamico deriva da quello standard per le miscele
compressibili e aggiunge ciò che serve al modello a due temperature: la
temperatura e l'energia vibro-elettronica come due campi in più, e tre termini
sorgente messi a disposizione del solver. Derivando da quello standard, per il
resto di OpenFOAM è un modello termodinamico come gli altri, e il caso lo
sceglie per nome come qualunque altro. È anche l'unico punto in cui si chiama
\mpp: il solver non parla mai direttamente con la libreria.

\subsection{\mpp}
\label{sec:mutation}

Questa sezione descrive la libreria che fornisce al progetto tutta la
termochimica e il modo in cui il progetto vi si collega.

\subsubsection{Che cos'è e a cosa serve}
\label{sec:mutation-cosa}

\mpp~\cite{scoggins2020} è una libreria scritta in C++, sviluppata
originariamente al von Karman Institute for Fluid Dynamics e distribuita come
codice aperto; il progetto ne usa la versione 1.1.3. È pensata proprio per
essere collegata a un codice CFD: non risolve equazioni di flusso, ma, dato lo
stato di un gas in un punto, ne calcola le proprietà che un solver non sa
calcolare da solo: termodinamiche, di trasporto, chimiche e di scambio di
energia fra i modi, dai flussi subsonici a quelli ipersonici, anche per gas
ionizzati. Il gas è descritto dalle specie che lo compongono, dal numero di
temperature, una oppure due come nel modello di Park, da una banca dati
termodinamica e dall'elenco delle reazioni. Lo stato si può dare in tre modi:
densità e temperature, densità ed energie, oppure frazioni in massa con
pressione e temperature. Da ognuno la libreria ricava tutto il resto, e dalle
energie trova da sé le temperature.

\subsubsection{Il collegamento con il progetto}
\label{sec:mutation-collegamento}

\mpp{} non fa parte di OpenFOAM: è compilata a parte, insieme al progetto, come
libreria dinamica, e il modello termodinamico ad alta entalpia descritto nella
sezione~\ref{sec:openfoam-moduli} e i programmi zero-dimensionali di
verifica ne chiamano le funzioni; il solver ci arriva attraverso il modello. All'avvio di una simulazione il
modello termodinamico legge nel caso la descrizione del gas -- specie, numero
di temperature, banca dati termodinamica, reazioni --, crea il gas
corrispondente in \mpp{} e associa ogni specie di OpenFOAM alla specie di
\mpp{} con lo stesso nome. Da quel momento, cella per cella e a ogni passo, gli
passa lo stato del gas e ne riceve le temperature e tutte le grandezze che
servono al solver. I dati del gas, cioè le costanti delle specie e le
reazioni, sono letti da una cartella del progetto che contiene i valori del
lavoro di riferimento.

\subsection{Il modello implementato}
\label{sec:implementato}

Questa sezione descrive che cosa fa il modello termodinamico ad alta entalpia:
con quali grandezze descrive il gas, quali termini sorgente calcola, come
aggira i limiti di \mpp{} e come il solver lo usa a ogni passo.

\subsubsection{Variabili e stato del gas}
\label{sec:implementato-stato}

In ogni cella il solver non segue le temperature, ma tre grandezze che si
conservano, come nel lavoro di riferimento: la densità di ciascuna specie,
l'energia del gas e la sua parte vibro-elettronica. Le temperature si ricavano
da queste. A ogni passo, dopo che il solver ha aggiornato le energie, il
modello le consegna a \mpp, che usa il modello a due temperature con la banca
dati RRHO e restituisce la temperatura traslazionale e quella
vibro-elettronica; con queste si aggiorna il legame fra pressione e densità.
All'avvio si fa il percorso inverso: dalle temperature assegnate nel caso,
\mpp{} calcola le energie iniziali di ogni cella. Nel passaggio va gestita una
differenza di convenzione: l'energia di OpenFOAM non comprende l'energia
chimica di formazione delle specie, quella di \mpp{} sì, e il modello la
aggiunge o la toglie a seconda della direzione. Il modello a polinomi del
template resta solo per costruire i campi che OpenFOAM si aspetta: tutte le
grandezze del gas vengono da \mpp, e la temperatura iniziale viene riletta dal
caso, perché il modello a polinomi l'avrebbe tagliata a 20\,000~K.

\subsubsection{I termini sorgente}
\label{sec:implementato-sorgenti}

Nel bagno termico le energie e la composizione cambiano solo per effetto dei
termini sorgente, calcolati cella per cella con le funzioni di \mpp. Il primo
è lo scambio fra vibrazione e traslazione, descritto dall'equazione di
Landau--Teller~\cite{landau1936}:
\[
  Q_{VT} = \sum_m \rho_m\,
  \frac{e_{ve,m}(\Ttr) - e_{ve,m}(\Tve)}{\tau_m},
\]
dove la somma è sulle molecole, $\rho_m$ è la densità della molecola $m$,
$e_{ve,m}$ la sua energia vibro-elettronica per unità di massa e $\tau_m$ il
suo tempo di rilassamento, calcolato da \mpp{} con la correlazione di Millikan
e White~\cite{millikan1963} e la correzione di Park. In parole: l'energia
vibrazionale si avvicina al valore che avrebbe alla temperatura
traslazionale, tanto più in fretta quanto più $\tau_m$ è piccolo. Il secondo è
la chimica: \mpp{} fornisce la produzione $\dot\omega_s$ di ogni specie, con le
costanti di Arrhenius del meccanismo valutate alla temperatura di Park
\[
  T_P = \Ttr^{\,a}\,\Tve^{\,1-a}, \qquad a = 0{,}7 .
\]
Le reazioni entrano poi nei bilanci di energia in due modi. Le molecole che si
dissociano portano via la loro energia vibro-elettronica e quelle che si
formano la aggiungono (modello non preferenziale),
\[
  Q_{CV} = \sum_s e_{ve,s}(\Tve)\,\dot\omega_s ;
\]
inoltre, poiché l'energia di OpenFOAM non comprende quella di formazione,
l'energia spesa per rompere i legami chimici va tolta a quella del gas con il
calore di reazione $Q_{ch} = -\sum_s h_{f,s}\,\dot\omega_s$, dove $h_{f,s}$ è
l'entalpia di formazione della specie $s$. In sintesi: le specie cambiano per
$\dot\omega_s$, l'energia vibro-elettronica per $Q_{VT} + Q_{CV}$ e l'energia
del gas, che la comprende, per $Q_{ch}$, così che l'energia totale, compresa
quella chimica, si conserva.

\subsubsection{Come sono stati aggirati i limiti di \mpp}
\label{sec:implementato-limiti}

I modelli di \mpp{} non coincidono in tutto con quelli del lavoro di
riferimento. Le differenze, e il modo in cui sono state trattate, sono queste:

\begin{itemize}[leftmargin=1.6em]
  \item \textbf{Energia che guida lo scambio V--T.} \mpp{} usa la sola
    energia vibrazionale, il lavoro di riferimento anche quella elettronica
    della molecola. Lo scambio è stato riscritto come nel lavoro di
    riferimento, con le energie e i tempi di rilassamento della libreria;
    senza livelli elettronici le due forme coincidono.
  \item \textbf{Temperatura delle dissociazioni.} \mpp{} usa sempre
    $\sqrt{\Ttr\Tve}$, cioè $a = 0{,}5$, e non si può cambiare. Per ottenere
    $a = 0{,}7$ si porta per un istante la libreria in uno stato in cui
    entrambe le temperature valgono $T_P$, si chiedono le velocità di reazione
    e si ripristina lo stato vero. Il procedimento vale solo per meccanismi
    fatti di sole dissociazioni, come quello dell'azoto usato qui.
  \item \textbf{Tempo di rilassamento nelle miscele.} \mpp{} media i tempi
    relativi ai vari partner di collisione in modo diverso dal lavoro di
    riferimento. Si è tenuta la versione della libreria, per non riscriverne
    il modello; l'effetto è discusso insieme ai risultati.
  \item \textbf{Una sola temperatura vibrazionale.} Nei casi risolti con il
    solver l'unica molecola presente è l'azoto, e il limite non conta. Il caso
    con azoto e ossigeno è risolto solo con un programma zero-dimensionale,
    che tiene separate le due vibrazioni e calcola a parte lo scambio V--V fra
    di esse.
  \item \textbf{Accoppiamento chimica-vibrazione.} \mpp{} offre solo il
    modello non preferenziale, ed è quello usato; il modello CVDV del lavoro
    di riferimento non è stato riprodotto.
  \item \textbf{Reazioni inverse.} Il set di reazioni per l'aria contiene due
    coppie di reazioni irreversibili una inversa dell'altra, che \mpp{} non
    accetta nello stesso meccanismo: sono state divise in due meccanismi,
    le cui velocità di reazione vengono sommate.
\end{itemize}

\subsubsection{Il solver}
\label{sec:implementato-solver}

Il solver a due temperature è derivato dal modulo density-based di OpenFOAM:
ne eredita il calcolo dei flussi, della densità, della quantità di moto e della
pressione, e riscrive la fase dell'energia. In questa fase, a ogni passo:

\begin{enumerate}[leftmargin=1.8em]
  \item risolve le equazioni delle specie, con la produzione chimica fornita
    da \mpp, e corregge le frazioni perché la loro somma resti uno;
  \item risolve l'equazione dell'energia del gas, uguale a quella del modulo
    originale più il calore di reazione;
  \item risolve l'equazione dell'energia vibro-elettronica, che per ora
    contiene solo la variazione nel tempo e i termini sorgente: manca il
    trasporto con il flusso, che nel bagno termico non serve;
  \item consegna le energie a \mpp, che restituisce le due temperature, e
    aggiorna il legame fra pressione e densità.
\end{enumerate}

\noindent I termini sorgente sono calcolati con lo stato del gas già noto e
trattati in modo esplicito. La fase dell'energia del modulo originale non
poteva essere riusata perché, subito dopo aver risolto l'equazione
dell'energia, ne ricava la temperatura: mancherebbero sia il calore di
reazione sia l'equazione vibro-elettronica, che deve essere risolta prima di
convertire le energie in temperature.

\subsubsection{Le equazioni risolte: forma generale e caso del bagno termico}
\label{sec:implementato-equazioni}

La Sezione~\ref{sec:implementato-sorgenti} ha dato i termini sorgente; questa
sezione riporta le equazioni di conservazione in cui entrano, prima nella
forma generale per cui \texttt{shockFluid} e \shockThermo{} sono scritti, valida
su una mesh qualunque, poi nella forma ridotta che viene effettivamente
integrata nei casi a bagno termico usati per la verifica
(Sezione~\ref{sec:programmi0D} e i casi \shockThermo{} corrispondenti).

Su una mesh generica, massa e quantità di moto sono risolte senza modifiche
rispetto al modulo originale \texttt{shockFluid},
\[
  \frac{\partial \rho}{\partial t} + \nabla\cdot(\rho\mathbf{U}) = 0 ,
  \qquad
  \frac{\partial (\rho\mathbf{U})}{\partial t}
  + \nabla\cdot(\rho\mathbf{U}\mathbf{U}) + \nabla p = \nabla\cdot\tau ,
\]
con lo sforzo viscoso newtoniano $\tau = \mu\left[\nabla\mathbf{U} +
(\nabla\mathbf{U})^{T} - \tfrac{2}{3}(\nabla\cdot\mathbf{U})\mathbf{I}\right]$,
dove $\mu$ è la viscosità dinamica. L'equazione dell'energia risolta da
\shockThermo{} aggiunge il calore di reazione a quella del modulo originale,
\[
  \frac{\partial (\rho e)}{\partial t} + \nabla\cdot(\boldsymbol\phi_p)
  + \frac{\partial (\rho K)}{\partial t}
  = Q_{ch} + \nabla\cdot(\tau\cdot\mathbf{U}) + \nabla\cdot(\kappa\nabla T) ,
\]
dove $K = \tfrac{1}{2}|\mathbf{U}|^2$ è l'energia cinetica per unità di massa,
$\kappa$ la conducibilità termica, e $\boldsymbol\phi_p$ il flusso convettivo
di energia attraverso ciascuna faccia, comprensivo del lavoro di pressione,
ricostruito con lo stesso schema di Kurganov--Tadmor usato per massa e
quantità di moto. Per analogia, una forma generale (non a bagno termico)
dell'equazione vibro-elettronica aggiungerebbe un termine convettivo e uno
diffusivo dello stesso tipo,
\[
  \frac{\partial (\rho\eve)}{\partial t} + \nabla\cdot(\phi\,\eve)
  = Q_{VT} + Q_{CV} + \nabla\cdot(\kappa_{ve}\nabla\Tve) ,
\]
ma, come notato nella Sezione~\ref{sec:implementato-solver}, sono
attualmente implementati solo la variazione nel tempo e i termini sorgente.

Nessuno dei termini di trasporto qui sopra è mai attraversato dai casi
risolti in questo progetto. Ogni caso è un bagno termico: una singola cella
di calcolo, con campo uniforme nello spazio, ogni condizione al contorno a
gradiente nullo e il gas fermo, $\mathbf{U}=0$.
In queste condizioni $\nabla\mathbf{U}=0$ e $\nabla T=0$ identicamente, in
ogni cella e a ogni passo temporale, quindi
\[
  \tau = 0, \qquad \kappa\nabla T = 0, \qquad K = 0, \qquad
  \boldsymbol\phi_p = 0 ,
\]
indipendentemente dal valore numerico di $\mu$ e $\kappa$; massa e quantità
di moto si riducono all'identità banale $0=0$. Quello che resta, ed è ciò
che viene effettivamente integrato, è
\[
  \frac{\partial (\rho e)}{\partial t} = Q_{ch} ,
  \qquad
  \frac{\partial (\rho \eve)}{\partial t} = Q_{VT} + Q_{CV} .
\]
Sono le stesse due equazioni differenziali ordinarie avanzate nel tempo dai
programmi 0D della Sezione~\ref{sec:programmi0D}, ed è questo che rende il
confronto fra le due implementazioni un vero controllo del solver, e non
della mesh o della discretizzazione.

\subsection{I programmi di verifica zero-dimensionali}
\label{sec:programmi0D}

Accanto al solver il progetto comprende quattro piccoli programmi che simulano
il bagno termico senza OpenFOAM, uno per ogni famiglia di casi del lavoro di
riferimento: l'azoto puro, la miscela di azoto molecolare e atomico, con o
senza reazioni, la miscela di azoto e ossigeno e l'aria a cinque specie.
Ognuno prepara il gas nello stato iniziale della figura da riprodurre, calcola
le energie contenute nel volume e avanza nel tempo con passi di un
nanosecondo, come il lavoro di riferimento: a ogni passo calcola i termini
sorgente, aggiorna l'energia vibro-elettronica e la composizione, e chiede a
\mpp{} le nuove temperature. L'energia totale, compresa quella chimica, resta
costante, perché il volume è chiuso e isolato. I termini sorgente sono
calcolati con le stesse funzioni usate dal solver, ed è questo che rende
significativo il confronto fra i due percorsi. Nei casi senza reazioni i
programmi calcolano anche la temperatura finale attesa, cercando la
temperatura unica alla quale il gas avrebbe la stessa energia dello stato
iniziale: è il valore verso cui le due temperature devono convergere.

Due programmi si discostano da questo schema. Quello per azoto e ossigeno tiene
separate le energie vibrazionali delle due molecole, ciascuna con la propria
temperatura, e aggiunge lo scambio V--V fra di esse con la formulazione di
Knab \textit{et al.}~\cite{knab1995} usata dal lavoro di riferimento e la
probabilità di scambio consigliata di $10^{-2}$; la sezione d'urto della coppia
azoto-ossigeno, che il lavoro di riferimento non riporta, è presa da
\hyfoam~\cite{hystrath}. Quello per l'aria usa invece una sola temperatura,
come fa il lavoro di riferimento per questo caso. I programmi sono anche molto
più rapidi del solver: il caso più lungo, un milione di passi, richiede circa
un secondo contro quasi mezz'ora.

\subsection{I dati del gas}
\label{sec:dati}

Le proprietà delle specie usate da \mpp{} sono quelle delle tabelle del lavoro
di riferimento: temperature caratteristiche di vibrazione, entalpie di
formazione e livelli elettronici delle cinque specie dell'aria. I parametri del
rilassamento vibrazionale sono quelli già contenuti in \mpp, tratti da
Park~\cite{park1993}, che per azoto e ossigeno coincidono con quelli del lavoro
di riferimento.

L'energia elettronica richiede una scelta. Il lavoro di riferimento la tratta
esplicitamente solo nel caso dell'azoto a 30\,000~K, dove confronta il gas con
e senza; per gli altri casi non lo dice, e il testo che commenta la miscela di
azoto molecolare e atomico parla anzi del contributo elettronico dell'azoto
atomico. Le temperature finali riportate nell'articolo, però, tornano solo
senza livelli elettronici: con i livelli elettronici, la temperatura finale
della miscela di azoto molecolare e atomico sarebbe di circa 18\,300~K invece
dei 24\,400~K dell'articolo, e quella della miscela di azoto e ossigeno di
circa 19\,200~K invece di 12\,100~K, mentre senza si ottengono 24\,353~K e
12\,133~K. Si sono quindi seguiti i risultati: il progetto ha due versioni dei
dati, una con i livelli elettronici, usata solo per il caso dell'azoto che li
include, e una con il solo livello fondamentale, usata in tutti gli altri
casi.

Le reazioni sono di due tipi. Per la dissociazione dell'azoto si usa la
reazione irreversibile $\mathrm{N_2 + N_2 \rightarrow 2\,N + N_2}$ con le
costanti di Park riportate nel lavoro di riferimento. Per l'aria a cinque
specie il lavoro di riferimento usa diciannove reazioni irreversibili, quindici
dissociazioni e quattro scambi, con costanti ricavate dalla teoria
Quantum-Kinetic~\cite{scanlon2015}, che cita ma non riporta: sono state prese
dai dati di \hyfoam~\cite{hystrath}. Per confronto è stato usato anche il
meccanismo per l'aria già presente in \mpp, con le costanti di Park.

% ---------------------------------------------------------------------------
% Traccia per le prossime sottosezioni (dalla vecchia sezione "Il modello"):
%  - decomposizione dell'energia nei modi (eq. 1-7 del paper) e le due T
%  - variabili conservate: rho_s, E, E_ve (eq. 23); perche' energie e non T
%  - equazione di stato: p = somma delle pressioni parziali (eq. 24)
%  - scambio V-T: Landau-Teller (eq. 8), tau di Millikan-White + Park (eq. 9-17)
%  - scambio V-V: Knab (eq. 18), solo nel caso N2-O2
%  - chimica: legge di azione di massa (eq. 27), Arrhenius (eq. 28),
%    temperatura di Park T^a Tv^(1-a) (eq. 29)
%  - accoppiamento chimica-vibrazione: modello non preferenziale (eq. 30-31)
%  - [fatto, vedi sec:implementato-equazioni] riepilogo del sistema risolto
%    nel bagno termico (eq. 22, 25, 26)
